\begin{table}[!ht]
\centering
\caption{Decomposição interseccional (raça $\times$ gênero) do gap de rendimento
\emph{vs.}\ o Homem Branco (grupo de referência). Dotações: parcela explicada por
características observáveis; Retornos: parcela não explicada (discriminação). Penalidade
extra: o quanto o gap da Mulher Negra excede a soma das penalidades de raça e de gênero
isoladas (efeito interseccional puro, Crenshaw, 1989). PNAD Contínua, população completa.}
\label{tab:interseccional}
\begin{tabular}{lrrrr}
\toprule
Grupo & Gap vs HB (\%) & Dotações (\%) & Retornos (\%) & Penal.\ extra (p.p.) \\
\midrule
Mulher Branca & 46,6 & −8,9 & 108,9 & --- \\
Homem Negro & 40,3 & 70,9 & 29,1 & --- \\
Mulher Negra & 96,4 & 29,6 & 70,4 & 9,5 \\
\bottomrule
\end{tabular}
\par\smallskip
\footnotesize\emph{Como ler:} ``Gap vs HB'' é a desvantagem de renda de cada grupo
\emph{vs.}\ o Homem Branco; Dotações $+$ Retornos $=100\%$ do gap. A Mulher Negra tem o
maior gap (96,4\%) e ainda uma penalidade \emph{extra} de 9,5 p.p.\ que não se reduz à
soma de ``ser negro'' e ``ser mulher'' --- a marca da interseccionalidade.
\end{table}
